% Pagina di prova ed esempio d'uso di tutti gli ambienti.
% Aggiungi un campione delle scritture e dei simboli del corso, compila, guarda il risultato.
% Rileggilo come modello, poi cancella il file e la sua riga \input.
% I contenuti sono solo esempi: non vanno riportati negli appunti.
\chapter{Prova tecnica}\label{cap:prova}

\begin{primadiiniziare}
\textbf{Di cosa parla:} il capitolo mostra tutti gli strumenti del template.

\textbf{Che cosa devi ricordare:} niente, è il primo capitolo.

\textbf{Alla fine saprai:} riconoscere ogni riquadro e ogni tipo di schema.
\end{primadiiniziare}

\begin{mappa}
\centering
\begin{tikzpicture}[node distance=1.8cm and .3cm]
\node[concetto=4.4cm] (c) {Strumenti degli appunti};
\node[dettaglio,below left=of c] (s) {Schemi};
\node[dettaglio,below=of c] (t) {Testo e citazioni};
\node[dettaglio,below right=of c] (n) {Numeri e codice};
\draw[freccia] (c) -- node[etichetta] {comprende} (s);
\draw[freccia] (c) -- node[etichetta] {comprende} (t);
\draw[freccia] (c) -- node[etichetta] {comprende} (n);
\end{tikzpicture}
\end{mappa}

La mappa ha un concetto centrale, «Strumenti degli appunti», che comprende tre parti: gli schemi, il testo con le citazioni, i numeri con il codice.

\section{Schemi}

\subsection{Mappa di un concetto}

Questa è una mappa concettuale. Ogni freccia dice quale relazione lega due concetti.

\begin{schema}[idea, sue parti ed esempio]
\centering
\begin{tikzpicture}[node distance=1.8cm and 1cm]
\node[concetto=3cm] (a) {Concetto};
\node[dettaglio=3cm,right=3.8cm of a] (b) {Definizione};
\node[esempio=3cm,below=of a] (d) {Esempio del corso};
\draw[freccia] (a) -- node[etichetta] {si definisce come} (b);
\draw[freccia] (a) -- node[etichetta] {si vede in} (d);
\end{tikzpicture}
\end{schema}

Lo schema dice che un concetto ha una definizione e si vede in un esempio del corso.

\subsection{Diagramma di flusso}

\begin{flusso}[come scegliere la forma di uno schema]
\centering
\begin{tikzpicture}[node distance=1.2cm and 1.6cm]
\node[estremo] (i) {Inizio};
\node[decisione,below=of i] (d) {Ci sono passi in sequenza?};
\node[processo=4.4cm,below=of d] (f) {Usa un diagramma di flusso};
\node[processo=4.4cm,right=of d] (m) {Usa una mappa concettuale};
\node[estremo,below=of f] (e) {Fine};
\draw[freccia] (i) -- (d);
\draw[freccia] (d) -- node[etichetta] {sì} (f);
\draw[freccia] (d) -- node[etichetta] {no} (m);
\draw[freccia] (f) -- (e);
\draw[freccia] (m) |- (e);
\end{tikzpicture}
\end{flusso}

Il diagramma fa una domanda: ci sono passi in sequenza? Se la risposta è sì, si usa un diagramma di flusso. Se la risposta è no, si usa una mappa concettuale.

\section{Testo e citazioni}
% nodiagramma: sezione dimostrativa senza concetti da schematizzare

Italiano con accenti: è, à, ì, ò, ù, perché, città. Parola inglese: \en{translation}\aut{Autore, Nome}\idx{concetto}. Titolo di un'opera: \emph{Titolo dell'opera}. Termine chiave: \textbf{termine}. Nome con pronuncia: Goethe \pron{ghöte}.

Greco: \gr{παρθένος}, \gr{Ἰδοὺ ἡ παρθένος ἐν γαστρὶ ἕξει}. Parole con legature evitate: fine, fluire, affitto, ufficio. Rimando: capitolo~\ref{cap:prova}.

\begin{cit}[english]
A quotation in the language of the source stays complete, however hard it may be to read.
\fontecit{(Fonte, anno)}
\end{cit}

\begin{parolesemplici}
La citazione resta intera. Questo riquadro la spiega con parole più semplici.
\end{parolesemplici}

\afronte[english]{The source text stands in this column, in its language, across several lines.}{La traduzione sta in questa colonna, su più righe.}

\begin{nonconfondere}
\textbf{Citazione} è il testo originale. \textbf{Parafrasi} è lo stesso testo detto con altre parole.
\end{nonconfondere}

\begin{permemoria}
Un trucco per ricordare va sempre dichiarato come tale.
\end{permemoria}

\begin{correzione}Prova di correzione.\end{correzione}
\begin{dubbio}Prova di dubbio.\end{dubbio}
\begin{dallibro}Prova dal libro (p.~1).\end{dallibro}
\soloappunti
Testo. \integrazione{prova di integrazione.}

\section{Numeri, tabelle e codice}
% nodiagramma: sezione dimostrativa senza concetti da schematizzare

\subsection{Una formula}

La somma dei primi $n$ numeri è:
\begin{equation}
S = \frac{n(n+1)}{2}
\end{equation}

\begin{simboli}
$S$ & la somma cercata & nessuna \\
$n$ & l'ultimo numero da sommare & nessuna \\
\end{simboli}

\begin{passoapasso}{la somma dei primi 4 numeri}
\begin{enumerate}
\item Sostituisco $n$ con 4: $S = \frac{4 \cdot (4+1)}{2}$.
\item Calcolo la parentesi, perché ha la precedenza: $4+1 = 5$, quindi $S = \frac{4 \cdot 5}{2}$.
\item Calcolo il prodotto: $4 \cdot 5 = 20$, quindi $S = \frac{20}{2}$.
\item Divido: $S = 10$. Controllo: $1+2+3+4 = 10$.
\end{enumerate}
\end{passoapasso}

\subsection{Tabelle e date}

\begin{center}
\begin{tabularx}{\linewidth}{@{}lYY@{}}
\toprule
& \textbf{Prima colonna} & \textbf{Seconda colonna} \\
\midrule
Riga & un testo abbastanza lungo da andare a capo dentro la cella della tabella & un altro testo abbastanza lungo da andare a capo \\
\bottomrule
\end{tabularx}
\end{center}

\begin{cronologia}
1813 & Primo evento di esempio \\
1923 & Secondo evento di esempio \\
\end{cronologia}

\subsection{Codice}

\begin{lstlisting}[language=C,numbers=left]
int main(void) { /* commento con è */
    return 0;
}
\end{lstlisting}

\begin{daricordare}
\begin{itemize}
\item Ogni capitolo ha la stessa struttura.
\item Gli schemi si aggiungono al testo.
\item Ogni formula ha la spiegazione dei simboli.
\end{itemize}
\end{daricordare}

\begin{insintesi}
Il template offre schemi, riquadri e tabelle pronte. Servono a dare più vie d'accesso agli stessi contenuti.
\end{insintesi}

\begin{ripasso}
\begin{enumerate}
\item Che cosa deve avere ogni freccia di una mappa concettuale?
\item Che cosa segue sempre una citazione difficile?
\end{enumerate}
\end{ripasso}

\begin{risposte}
\begin{enumerate}
\item Un'etichetta che dice la relazione (sezione «Schemi»).
\item Un riquadro «In parole semplici» (sezione «Testo e citazioni»).
\end{enumerate}
\end{risposte}
